% style_macros.tex — simple LaTeX macros for the S.U.R.E. paper (project_000).
% % style_macros.tex — simple LaTeX macros for the S.U.R.E. paper (project_000).
% % style_macros.tex — simple LaTeX macros for the S.U.R.E. paper (project_000).
% % style_macros.tex — simple LaTeX macros for the S.U.R.E. paper (project_000).
% \input{style_macros} from the preamble of final_paper.tex.
% Provides \question{} and \takeaway{} environments plus a \formalization{} environment used for
% the 4 Key Formalizations (explicitly NOT labeled as theorems — this is an empirical-only paper,
% per track_decomposition.json's recommended_track="empirical").

\usepackage{amsmath}
\usepackage{amsthm}
\usepackage{xcolor}
\usepackage{tcolorbox}
\tcbuselibrary{skins,breakable}

% ---- Formalization environment (not "Theorem") ----
\newtheoremstyle{formalizationstyle}
  {6pt}{6pt}{\itshape}{}{\bfseries}{.}{.5em}{}
\theoremstyle{formalizationstyle}
\newtheorem{formalization}{Formalization}

% ---- \question{} — used sparingly in Introduction/Discussion to name a research question ----
\newcommand{\question}[1]{%
  \par\medskip\noindent\textbf{Question.}\ #1\par\medskip
}

% ---- \takeaway{} — used in Introduction to state an explicit, section-pointing takeaway ----
\newtcolorbox{takeawaybox}{%
  colback=gray!5, colframe=gray!40, boxrule=0.4pt, arc=1pt,
  left=6pt, right=6pt, top=4pt, bottom=4pt, breakable
}
\newcommand{\takeaway}[1]{%
  \begin{takeawaybox}\textbf{Takeaway.}\ #1\end{takeawaybox}
}

% ---- \sourceclaim{} — inline marker for citing the exact raw-artifact source of a number ----
% Usage: state the number, then \sourceclaim{EXP07/exp07_conf_curve_summary.json}
\newcommand{\sourceclaim}[1]{\text{\footnotesize\color{gray}[#1]}}

% ---- \disclosure{} — visually sets off a must-disclose-early callout paragraph (e.g. the
% 0.719-vs-0.782 recall distinction) as a named callout, not a figure caption ----
\newtcolorbox{disclosurebox}{%
  colback=yellow!5, colframe=black!50, boxrule=0.5pt, arc=1pt,
  left=6pt, right=6pt, top=4pt, bottom=4pt, breakable,
  title={\small\bfseries Disclosure}, fonttitle=\small, coltitle=black,
  colbacktitle=gray!15
}
\newcommand{\disclosure}[1]{\begin{disclosurebox}#1\end{disclosurebox}}
 from the preamble of final_paper.tex.
% Provides \question{} and \takeaway{} environments plus a \formalization{} environment used for
% the 4 Key Formalizations (explicitly NOT labeled as theorems — this is an empirical-only paper,
% per track_decomposition.json's recommended_track="empirical").

\usepackage{amsmath}
\usepackage{amsthm}
\usepackage{xcolor}
\usepackage{tcolorbox}
\tcbuselibrary{skins,breakable}

% ---- Formalization environment (not "Theorem") ----
\newtheoremstyle{formalizationstyle}
  {6pt}{6pt}{\itshape}{}{\bfseries}{.}{.5em}{}
\theoremstyle{formalizationstyle}
\newtheorem{formalization}{Formalization}

% ---- \question{} — used sparingly in Introduction/Discussion to name a research question ----
\newcommand{\question}[1]{%
  \par\medskip\noindent\textbf{Question.}\ #1\par\medskip
}

% ---- \takeaway{} — used in Introduction to state an explicit, section-pointing takeaway ----
\newtcolorbox{takeawaybox}{%
  colback=gray!5, colframe=gray!40, boxrule=0.4pt, arc=1pt,
  left=6pt, right=6pt, top=4pt, bottom=4pt, breakable
}
\newcommand{\takeaway}[1]{%
  \begin{takeawaybox}\textbf{Takeaway.}\ #1\end{takeawaybox}
}

% ---- \sourceclaim{} — inline marker for citing the exact raw-artifact source of a number ----
% Usage: state the number, then \sourceclaim{EXP07/exp07_conf_curve_summary.json}
\newcommand{\sourceclaim}[1]{\text{\footnotesize\color{gray}[#1]}}

% ---- \disclosure{} — visually sets off a must-disclose-early callout paragraph (e.g. the
% 0.719-vs-0.782 recall distinction) as a named callout, not a figure caption ----
\newtcolorbox{disclosurebox}{%
  colback=yellow!5, colframe=black!50, boxrule=0.5pt, arc=1pt,
  left=6pt, right=6pt, top=4pt, bottom=4pt, breakable,
  title={\small\bfseries Disclosure}, fonttitle=\small, coltitle=black,
  colbacktitle=gray!15
}
\newcommand{\disclosure}[1]{\begin{disclosurebox}#1\end{disclosurebox}}
 from the preamble of final_paper.tex.
% Provides \question{} and \takeaway{} environments plus a \formalization{} environment used for
% the 4 Key Formalizations (explicitly NOT labeled as theorems — this is an empirical-only paper,
% per track_decomposition.json's recommended_track="empirical").

\usepackage{amsmath}
\usepackage{amsthm}
\usepackage{xcolor}
\usepackage{tcolorbox}
\tcbuselibrary{skins,breakable}

% ---- Formalization environment (not "Theorem") ----
\newtheoremstyle{formalizationstyle}
  {6pt}{6pt}{\itshape}{}{\bfseries}{.}{.5em}{}
\theoremstyle{formalizationstyle}
\newtheorem{formalization}{Formalization}

% ---- \question{} — used sparingly in Introduction/Discussion to name a research question ----
\newcommand{\question}[1]{%
  \par\medskip\noindent\textbf{Question.}\ #1\par\medskip
}

% ---- \takeaway{} — used in Introduction to state an explicit, section-pointing takeaway ----
\newtcolorbox{takeawaybox}{%
  colback=gray!5, colframe=gray!40, boxrule=0.4pt, arc=1pt,
  left=6pt, right=6pt, top=4pt, bottom=4pt, breakable
}
\newcommand{\takeaway}[1]{%
  \begin{takeawaybox}\textbf{Takeaway.}\ #1\end{takeawaybox}
}

% ---- \sourceclaim{} — inline marker for citing the exact raw-artifact source of a number ----
% Usage: state the number, then \sourceclaim{EXP07/exp07_conf_curve_summary.json}
\newcommand{\sourceclaim}[1]{\text{\footnotesize\color{gray}[#1]}}

% ---- \disclosure{} — visually sets off a must-disclose-early callout paragraph (e.g. the
% 0.719-vs-0.782 recall distinction) as a named callout, not a figure caption ----
\newtcolorbox{disclosurebox}{%
  colback=yellow!5, colframe=black!50, boxrule=0.5pt, arc=1pt,
  left=6pt, right=6pt, top=4pt, bottom=4pt, breakable,
  title={\small\bfseries Disclosure}, fonttitle=\small, coltitle=black,
  colbacktitle=gray!15
}
\newcommand{\disclosure}[1]{\begin{disclosurebox}#1\end{disclosurebox}}
 from the preamble of final_paper.tex.
% Provides \question{} and \takeaway{} environments plus a \formalization{} environment used for
% the 4 Key Formalizations (explicitly NOT labeled as theorems — this is an empirical-only paper,
% per track_decomposition.json's recommended_track="empirical").

\usepackage{amsmath}
\usepackage{amsthm}
\usepackage{xcolor}
\usepackage{tcolorbox}
\tcbuselibrary{skins,breakable}

% ---- Formalization environment (not "Theorem") ----
\newtheoremstyle{formalizationstyle}
  {6pt}{6pt}{\itshape}{}{\bfseries}{.}{.5em}{}
\theoremstyle{formalizationstyle}
\newtheorem{formalization}{Formalization}

% ---- \question{} — used sparingly in Introduction/Discussion to name a research question ----
\newcommand{\question}[1]{%
  \par\medskip\noindent\textbf{Question.}\ #1\par\medskip
}

% ---- \takeaway{} — used in Introduction to state an explicit, section-pointing takeaway ----
\newtcolorbox{takeawaybox}{%
  colback=gray!5, colframe=gray!40, boxrule=0.4pt, arc=1pt,
  left=6pt, right=6pt, top=4pt, bottom=4pt, breakable
}
\newcommand{\takeaway}[1]{%
  \begin{takeawaybox}\textbf{Takeaway.}\ #1\end{takeawaybox}
}

% ---- \sourceclaim{} — inline marker for citing the exact raw-artifact source of a number ----
% Usage: state the number, then \sourceclaim{EXP07/exp07_conf_curve_summary.json}
\newcommand{\sourceclaim}[1]{\text{\footnotesize\color{gray}[#1]}}

% ---- \disclosure{} — visually sets off a must-disclose-early callout paragraph (e.g. the
% 0.719-vs-0.782 recall distinction) as a named callout, not a figure caption ----
\newtcolorbox{disclosurebox}{%
  colback=yellow!5, colframe=black!50, boxrule=0.5pt, arc=1pt,
  left=6pt, right=6pt, top=4pt, bottom=4pt, breakable,
  title={\small\bfseries Disclosure}, fonttitle=\small, coltitle=black,
  colbacktitle=gray!15
}
\newcommand{\disclosure}[1]{\begin{disclosurebox}#1\end{disclosurebox}}
